% !TEX root = ../main.tex
%% Corrigé du sujet type bac n° 2 — résultats vérifiés par programme (bac_s2.py)

\corrigeSujet{2}

\begin{corrigeBac}{4}{Questionnaire à choix multiples justifié}
\begin{questions}
  \item \textbf{Réponse b).} Pour $x\neq0$ : $\dfrac{e^{x}-x^{2}}{2e^{x}+1}=\dfrac{1-x^{2}e^{-x}}{2+e^{-x}}$, et $x^{2}e^{-x}\to0$
    par croissances comparées : la limite vaut $\frac12$. \pts{1 pt}
  \item \textbf{Réponse d).} En posant $u(x)=x^{2}$, on a $f=\frac12u'e^{u}$, de primitive $\frac12e^{u}$. Vérification :
    $\left(\frac12e^{x^{2}}\right)'=\frac12\times2x\,e^{x^{2}}=x\,e^{x^{2}}$. \pts{0,75 pt}
  \item \textbf{Réponse a).} Sur $\intOO1{+\infty}$, l'équation équivaut à $\ln\big(x(x-1)\big)=\ln6$, soit $x^{2}-x-6=0$, de
    solutions $-2$ et $3$ ; seule $3$ appartient à l'intervalle. \pts{0,75 pt}
  \item \textbf{Réponse c).} On choisit $2$ as parmi $4$, puis $3$ autres cartes parmi les $28$ qui ne sont pas des as :
    $\binom42\times\binom{28}{3}=6\times3\,276=19\,656$ mains. \pts{0,75 pt}
  \item \textbf{Réponse d).} $P(X\geq1)=1-P(X=0)=1-0{,}7^{10}\approx0{,}972$. \pts{0,75 pt}
\end{questions}
\end{corrigeBac}

\begin{corrigeBac}{5}{Suite arithmético-géométrique}
\begin{questions}
  \item $u_1=0{,}9\times2\,000+300=\res{2\,100}$ et $u_2=0{,}9\times2\,100+300=\res{2\,190}$. \pts{0,5 pt}
  \item Chaque mois, il reste $90\,\%$ des abonnés, soit $0{,}9\,u_n$, auxquels s'ajoutent $300$ nouveaux abonnés :
    $u_{n+1}=0{,}9\,u_n+300$. \pts{0,5 pt}
  \item a) $v_{n+1}=u_{n+1}-3\,000=0{,}9\,u_n-2\,700=0{,}9\,(u_n-3\,000)=0{,}9\,v_n$ : la suite $(v_n)$ est géométrique de raison
    $\res{0{,}9}$ et de premier terme $v_0=\res{-1\,000}$.\\
    b) On en déduit $v_n=-1\,000\times0{,}9^{n}$, puis $u_n=v_n+3\,000=3\,000-1\,000\times0{,}9^{n}$. \pts{1,5 pt}
  \item $u_{n+1}-u_n=-1\,000\times0{,}9^{n}(0{,}9-1)=100\times0{,}9^{n}>0$ : la suite est $\res{\text{croissante}}$. Comme
    $0<0{,}9<1$, $\lim0{,}9^{n}=0$ et $\res{\lim u_n=3\,000}$ : le nombre d'abonnés augmente et se stabilise autour de $3\,000$. \pts{1 pt}
  \item a) $u_n>2\,900\iff0{,}9^{n}<0{,}1\iff n\ln0{,}9<\ln0{,}1\iff n>\dfrac{\ln0{,}1}{\ln0{,}9}\approx21{,}85$ (car $\ln0{,}9<0$).
    Le plus petit entier convenable est $\res{n=22}$ : l'objectif est atteint au bout de $22$ mois, soit au
    1\textsuperscript{er} novembre 2027.\\
    b) La boucle tourne tant que le seuil n'est pas dépassé :
\begin{codecote}[0.5]
\begin{pythoncol}
def seuil(s):
    u = 2000
    n = 0
    while u <= s:
        u = 0.9 * u + 300
        n = n + 1
    return n
\end{pythoncol}
\cote
\code{seuil(2900)} renvoie \code{22}. On a $u_{21}\approx2\,890{,}6$ et $u_{22}\approx2\,901{,}5$.
\end{codecote}
    \pts{1,5 pt}
\end{questions}
\end{corrigeBac}

\begin{corrigeBac}{4}{Géométrie dans l'espace}
\begin{questions}
  \item Les vecteurs normaux $\vec{n_1}\,(1\pv1\pv-1)$ et $\vec{n_2}\,(2\pv-1\pv1)$ ne sont pas colinéaires (leurs deux premières
    coordonnées ne sont pas proportionnelles) : les plans sont $\res{\text{sécants}}$. \pts{0,5 pt}
  \item Pour tout réel $t$, le point $(1\pv t\pv t)$ vérifie $1+t-t-1=0$ et $2-t+t-2=0$ : cette droite est contenue dans
    $\mathcal{P}_1\cap\mathcal{P}_2$, qui est une droite. C'est donc $\Delta$. \pts{0,75 pt}
  \item $\vec{n_1}\cdot\vec{n_2}=2-1-1=0$ : les plans sont $\res{\text{perpendiculaires}}$. \pts{0,5 pt}
  \item $\Delta$ est dirigée par $\vec u\,(0\pv1\pv1)$, qui est donc normal à $\mathcal{Q}$ : $y+z+d=0$, avec $1+3+d=0$, soit
    $\res{\mathcal{Q}:\ y+z-4=0}$. \pts{0,5 pt}
  \item $t+t-4=0$ donne $t=2$ : $\res{H(1\pv2\pv2)}$. \pts{0,5 pt}
  \item $H\in\Delta$ et $\vect{MH}$ est orthogonal à $\Delta$, puisque $M$ et $H$ appartiennent au plan $\mathcal{Q}$ orthogonal à $\Delta$ :
    $H$ est le projeté orthogonal de $M$ sur $\Delta$. Avec $\vect{MH}\,(1\pv1\pv-1)$, la distance vaut $MH=\res{\sqrt3}$. \pts{0,5 pt}
  \item $H$ est le milieu de $[MM']$ : $M'=2H-M$, soit $\res{M'(2\pv3\pv1)}$. \pts{0,75 pt}
\end{questions}
\end{corrigeBac}

\begin{corrigeBac}{5}{Fonction logarithme et intégration}
\partieExo{A}{Étude de la fonction}
\begin{questions}
  \item En $0$ : $x^{2}\to0$ et $-2\ln x\to+\infty$, donc $\res{\lim\limits_{x\to0}f(x)=+\infty}$. En $+\infty$ :
    $f(x)=x^{2}\left(1-2\frac{\ln x}{x^{2}}\right)$ et $\frac{\ln x}{x^{2}}\to0$, donc $\res{\lim\limits_{x\to+\infty}f(x)=+\infty}$. \pts{1 pt}
  \item $f'(x)=2x-\frac2x=\frac{2(x^{2}-1)}{x}=\frac{2(x-1)(x+1)}{x}$, du signe de $x-1$ sur $\intOO0{+\infty}$.
\begin{center}
\begin{tableau}
  \otTab[2]{$x$/0.8,$f'(x)$/0.8,$f$/1.6}{$0$,$1$,$+\infty$}
  \tkzTabLine{d,-,z,+,}
  \tkzTabVar{D+/$+\infty$,-/$1$,+/$+\infty$}
\end{tableau}
\end{center}
    \pts{1 pt}
  \item Le minimum de $f$ sur $\intOO0{+\infty}$ est $f(1)=1$ : pour tout $x>0$, $f(x)\geq1$. \pts{0,25 pt}
  \item $f''(x)=2+\frac2{x^{2}}>0$ : $f$ est $\res{\text{convexe}}$. La courbe est au-dessus de ses tangentes, en particulier de la
    tangente en $1$, d'équation $y=1$ (car $f'(1)=0$) : on retrouve $f(x)\geq1$. \pts{0,75 pt}
\end{questions}
\partieExo{B}{Équation et aire}
\begin{questions}[resume]
  \item Sur $\intOF01$, $f$ est continue et strictement décroissante de $+\infty$ à $1$, et $2>1$ : l'équation $f(x)=2$ y a une unique
    solution $a\approx0{,}40$. De même sur $\intFO1{+\infty}$, où $f$ croît strictement de $1$ à $+\infty$ : une unique solution $b$, et
    $f(1{,}77)\approx1{,}991<2<f(1{,}78)\approx2{,}015$ donne $\res{1{,}77<b<1{,}78}$. \pts{1 pt}
  \item $\int_1^{e}f(x)\,\mathrm{d}x=\left[\frac{x^{3}}{3}\right]_1^{e}-2\Big[x\ln x-x\Big]_1^{e}=\frac{e^{3}-1}{3}-2
    =\res{\tfrac{e^{3}-7}{3}}\approx4{,}36$ : comme $f>0$, c'est l'aire (en u.a.) du domaine situé sous $\mathcal{C}_f$, entre $x=1$ et
    $x=e$. \pts{1 pt}
\end{questions}
\end{corrigeBac}
